%%%% This is -*- mode: latex; mode: flyspell -*-

\documentclass{article}

%% \usepackage[english]{babel}
%% \usepackage[latin1]{inputenc}
%% \usepackage[T1]{fontenc}

\usepackage[urw-garamond]{mathdesign}
%\renewcommand{\ttdefault}{cmtt}
%\usepackage{newcent}
%\usepackage{mathpazo}

\setcounter{secnumdepth}{5}
\setcounter{tocdepth}{5}

\usepackage{geometry}
\usepackage{verbatim}
\usepackage{cite}
\usepackage{listings}

\lstdefinelanguage{siexml}
  {morekeywords={xml,sie,version,xmlns,decoder,loop,var,bits,type,endian,
                 value,group,tag,id,format,read,sample,octets,encoding,
                 ch,test,index,dim,set,data,name},
   sensitive=false,
   morecomment=[s]{<!--}{-->},
   morestring=[b]",
   showstringspaces=false,
   columns=fullflexible
  }

\lstdefinelanguage{rng}{
  morekeywords={attribute,div,element,
    empty,external,grammar,include,inherit,list,
    mixed,notAllowed,parent,string,
    text,token},
  sensitive=false,
  morecomment=[l]{\#},
  morestring=[b]",
  columns=fullflexible
}

\lstnewenvironment{grammar}
  {\vbox\bgroup\lstset{language=rng,frame=simple}}
  {\egroup}

\lstnewenvironment{siexmlexample}
  {\vbox\bgroup\lstset{language=siexml,frame=simple}}
  {\egroup}

\newcommand{\relaxng}[0]{\hbox{RELAX~NG}}
\newcommand{\xml}[1]{\texttt{#1}}
\newcommand{\var}[1]{\texttt{#1}}
\newcommand{\lit}[1]{\texttt{#1}}

\newcommand{\fixme}{}
\newcommand{\incomplete}{\emph{(\ldots incomplete\ldots)}}

\newcommand{\undef}{\emph{undefined}}
\newcommand{\libsie}{\xml{libsie}}


\usepackage{multicol}

\newcommand{\siever}{1.0}

\title{The SIE Format}

\author{
Brian Downing \\ HBM, Inc. \and
Chris LaReau \\ HBM, Inc.
}

\date{December, 2009}

\begin{document}

\maketitle

\begin{abstract}
  SIE is a new data transmission and storage format for engineering
  data.  SIE is designed to be flexible, self-describing, streamable,
  and robust, overcoming some of the limitations of existing
  engineering data formats.  This is accomplished with a simple block
  structure and an XML description of how to read both the block
  structure and contained binary data.  It is also very simple to
  write, as some data loggers have very limited processor and memory
  capacity.  SIE has been successfully implemented on SoMat's eDAQ
  line of data acquisition systems, and a library for reading SIE
  files has been implemented in portable C.
\end{abstract}

\tableofcontents

\section{Introduction}

\subsection{Requirements}

Many of the requirements of a data format are dictated by the end-user
of the format --- those who read it.  For instance, at the time of
this writing, an 32 gigabyte CompactFlash card is available for
around US\$80.00, and the trend of increasing capacity and decreasing
price continues at an impressive rate.  Putting 16 or more
gigabytes of storage into a small data acquisition system is going to
be increasingly common.  Users of data acquisition systems and other
producers of engineering data want to be able to use all of this
space.

For situations involving long term testing, it is desired to be able
to read some or all of the data while testing is in progress.  In other
words, data stored in a file should be readable while the file is still
being written to.

As more information enters the digital world, it is becoming more
realistic to want many kinds of information associated with
engineering data.  A data format should be flexible enough to be able
to store arbitrary text, pictures, audio, video, and other information
not yet conceived.

Finally, and most importantly, nobody wants to lose data they've
collected.  Testing is very expensive, and hard or impossible to
repeat.  Any data storage format needs to be as robust as possible to
reduce the frequency of problems, and to make recovering data as easy
as possible when catastrophe occurs.

Additionally, some requirements are driven by implementers of hardware
and software that writes engineering data.  When collecting data, it
usually comes packed in some specific binary format.  An ideal data
format would offer flexible binary storage options so that the
incoming data could be directly written out, or at the very least with
minimum changes.

Also with the above, a good format should allow the details of the
binary data stored to be changed, without changing the engineering
values represented, and without requiring modifications to the reading
software.

Finally, the format should be very easy to write.  Ideally, it should
be feasible to write the format on microcontrollers with only
kilobytes of RAM.

\subsection{Problems with existing formats}

One of the primary reasons for creating a new data format is the
presence of many problems in existing formats.  Almost all existing
engineering formats have certain undesirable traits.

For instance, most existing engineering data formats are not
\emph{streamable} --- that is, they require random access to write and
read, usually because they have forward-pointing intra-file
references.  This limits the ability to send data in real time across
a network.  Also, the presence of forward pointers makes a file format
more complicated and error-prone.  If a file is corrupted and key
linking information is destroyed, it can be difficult or impossible to
recover otherwise unaffected data.

In some cases, files may not be in a consistent, readable state
until writing has stopped, either due to missing forward-pointing
offsets or other design issues.  This is particularly awkward for a
data acquisition system, as it is desirable to be able to fetch or
monitor data while acquisition is in progress.

In addition, many engineering data formats use 32-bit offsets
internally.  This limits files to no more than $2^{31}-1$ or $2^{32}$
bytes, depending on whether the offsets are interpreted as signed or
unsigned.  This has become a constraint as small, cheap flash devices
larger than four gigabytes have become available.

Finally, many existing data formats have limited metadata
representations.  Some have only fixed-field metadata; they may have
fields like channel name and sample rate, but no way to add other
pieces of metadata.  Other formats have an extensible metadata format
but do not use it pervasively or consistently.

\subsection{Design goals}

The following outline describes the major design goals for SIE:

\begin{itemize}
\item Flexible
  \begin{itemize}
  \item It should be possible to encode a wide variety of internal
    data formats without needing to extend the SIE specification.
  \end{itemize}
\item Conceptually simple and self-describing
  \begin{itemize}
  \item It should be possible to look at an SIE file and write code
    to extract data from it.
  \item Explicit data model - no implicit knowledge.
  \end{itemize}
\item No arbitrary limits:
  \begin{itemize}
  \item File size.
  \item Data formats.
  \item Metadata size, format, or language.
  \end{itemize}
\end{itemize}

\begin{itemize}
\item Streamable
  \begin{itemize}
  \item No forward data pointers.
  \item Append-only file writing.
  \end{itemize}
\item Safe
  \begin{itemize}
  \item Transmission or storage error detection.
  \item Maximal recoverability in case of corruption.
  \end{itemize}
\item Efficient
  \begin{itemize}
  \item Size:  minimal overhead as a percentage of data size for
    real-world examples.
  \item Processing:  should be possible to store data in native
    transducer format.
  \end{itemize}
\end{itemize}

\section{Design}

At its core, SIE is a block-structured container for data.  It defines
an XML \cite{xml} schema for describing both the structure of the data
and their associated metadata.  It also defines an optional index
format to speed up file access.

\subsection{Data model}
\label{data-model-sec}

SIE presents a flexible, unified data model to the end user.  Data are
presented as a table of rows and columns.  For example, a time series
channel could be represented as the following:

\begin{center}
  \begin{tabular}{r|r@{.}lr@{.}l}
    & \multicolumn{2}{l}{\textbf{Dimension 0}} & 
      \multicolumn{2}{l}{\textbf{Dimension 1}} \\ \hline
    \textbf{Row 0} & 0&0 &  0&42 \\
    \textbf{Row 1} & 0&1 & -0&20 \\
    \textbf{Row 2} & 0&2 &  0&13 \\
    \textbf{Row 3} & 0&3 &  0&06 \\
    \textbf{Row 4} & 0&4 & -0&23 \\
    $\cdots$ & \multicolumn{2}{l}{$\cdots$} & \multicolumn{2}{l}{$\cdots$}
  \end{tabular}
\end{center}

A dimension (column) may contain either numerical values or ``raw''
binary data.

The metadata for this example would contain additional information,
such as the fact that dimension 0 is \emph{time} and the applicable
unit is \emph{seconds}, while dimension 1 is the actual data
measurement, for which the unit is \emph{millivolts}.

Ultimately, all SIE data are presented in this basic form.  It is
notable that, if given an SIE file with completely unknown data, all
of the data values will be immediately visible even if the details of
the intended representation are unknown.

Metadata are represented as \emph{tags}, which are pairings of a
textual name to arbitrary binary data.  Tags can exist an any level of
the SIE hierarchy, as will be detailed below.  The metadata for the
example above could be represented as:

\begin{lstlisting}[language=siexml]
<ch id="0" name="example">
  <dim index="0">
    <tag id="core:label">time</tag>
    <tag id="core:units">seconds</tag>
  </dim>
  <dim index="1">
    <tag id="core:label">measurement</tag>
    <tag id="core:units">millivolts</tag>
  </dim>
</ch>
\end{lstlisting}

To help standardize and document the many ways data and metadata can
be represented on this basic model, SIE has \emph{schemas}, which are
documented sets of data and metadata representation standards.
Schemas are identified by a namespacing system --- for example, all
tags beginning with \xml{core:} belong to the \xml{core} schema.  The
\xml{sie} schema is reserved for items defined in the SIE format
proper, and will be described later in this document.  Two other
schemas exist at the time of this writing:  the aforementioned
\xml{core}, which describes what is believed to be metadata that will
be applicable to most engineering data, and \xml{somat}, which
describes metadata and data representations specific to SoMat's line
of data acquisition systems.  Please see the \xml{core} and
\xml{somat} schema documents for more information.

Note that data can be \emph{stored} in SIE in almost any conceivable
format.  Most of the SIE format is in fact a mechanism to convert
arbitrary binary data in a block-structured container into the unified
data model described above.

\subsection{Block structure}
\label{block-structure-sec}

SIE is output as a stream of bytes, which is composed of a series of
data blocks in sequence without intervening padding.  This is depicted
in figure~\ref{data-stream-fig}.  A block is composed of several
parts, namely a \emph{size}, \emph{group}, \emph{sync word},
\emph{payload}, and \emph{checksum}.  The size is repeated on both
ends of the block.  All of these except the payload are 32-bit
unsigned integers in network (big-endian) byte order.

\def\blockline{
  size & group & sync & 
  \hspace{1cm}payload\hspace{1cm} 
  & check & size \\ \hline
}
\def\blockstructure{
  \begin{center}
    \begin{tabular}{|c|c|c|c|c|c|}    \hline
      \blockline
      \blockline
      \multicolumn{6}{c}{$\cdots$} \\ \hline
      \blockline
    \end{tabular}
  \end{center}
}

\begin{figure}
  \blockstructure
  \caption{An SIE data stream}  \label{data-stream-fig}
\end{figure}

The \emph{size} is simply the size of the block in bytes, including
both size fields.  Having the size on both sides of the block allows
SIE data in a file to be traversed both forwards and backwards, and
also provides one level of consistency checking --- if the sizes don't
agree, something is wrong.

The \emph{group} identifies what data are contained in the
\emph{payload} of the block.  Groups 0 and 1 are always the XML
metadata and index blocks respectively.  The purpose of all other
groups is defined in the XML metadata.

The \emph{sync word} is always the hexadecimal value
\lit{51EDA7A0}.  This looks vaguely like ``SIEDATA0.''  The sync
word is used to find the beginning of a block (or ``resync'') in cases
such as damaged files or joining a broadcast stream in the middle.
Also, as it has some high bits set, if an SIE file is transmitted over
a medium that is not 8-bit clean, the sync word will be damaged,
quickly indicating corruption even if checksums are not present.

The \emph{payload} is the actual data contained in the block.  It is
interpreted based on the \emph{group} of the block.  As can be
deduced, the length of the payload is always $\emph{size} - 20$.

Finally, the \emph{checksum} is a CRC32 checksum of all bytes from the
first byte of the first \emph{size} through the last byte of the
\emph{payload}.  The checksum is optional, as some devices may not
have the processor capacity to compute it, but if it is not present,
the checksum field must be set to zero.  When reading a block, a zero
checksum is always valid.

A block with a payload size of zero indicates that no more blocks of
that group will occur.  This is intended to be used to indicate end of
data in real-time streaming scenarios.

\def\blocktable{
  \begin{center}
    \begin{tabular}{rrll}
      \textbf{Offset} & \textbf{Size} & 
        \textbf{Name} & \textbf{Description} \\ \hline
    $0$      & $4$ & \emph{size}      & Total size of the block in bytes. \\ 
    $4$      & $4$ & \emph{group}     & Group identifier for this block.  \\ 
    $8$      & $4$ & \emph{sync word} & Constant \lit{51EDA7A0}.          \\ 
    $12$     & $n$ & \emph{payload}   & $n$ bytes of data.                \\ 
    $n + 12$ & $4$ & \emph{checksum}  &
        CRC32 of bytes $0$ through $n + 11$ inclusive.                    \\ 
    $n + 16$ & $4$ & \emph{size 2}    & Same as \emph{size}.              \\ 
    \multicolumn{4}{c}{\emph{(All sizes and offsets in bytes.)}}
    \end{tabular}
  \end{center}
}

\begin{table}
  \blocktable
  \caption{An SIE block}  \label{block-tab}
\end{table}

\subsection{Predefined groups}

There are only two groups defined in the core SIE standard:  group 0
blocks contain the SIE XML metadata, and group 1 blocks are index
blocks.

\subsubsection{XML metadata (group 0)}

The SIE XML metadata defines how to interpret the SIE data.  It
contains \emph{tags}, which are generalized metadata relating an
arbitrary text key to a completely arbitrary value; \emph{channels},
which are groupings of engineering data (composed of multiple
\emph{dimensions}) with their associated metadata; and \emph{tests},
which are groupings of channels.  Tags can appear at all levels of the
hierarchy.

Additionally, the SIE XML metadata contains descriptions of all binary
formats contained within the file, including the block structure
defined above.  These descriptions are contained within
\emph{decoders}, which are small programs to read and interpret binary
data.

The XML representation of the SIE metadata is defined as the
sequential concatenation of the payload of every block with group ID
0.  The standard SIE XML metadata preamble is:

\begin{lstlisting}[language=siexml,frame=simple]
<?xml version="1.0" encoding="UTF-8"?>
<sie version="1.0" xmlns="http://www.somat.com/SIE">
<!-- SIE format standard definitions: -->
 <!-- SIE stream decoder: -->
 <decoder id="0">
  <loop>
   <read var="size" bits="32" type="uint" endian="big"/>
   <read var="group" bits="32" type="uint" endian="big"/>
   <read var="syncword" bits="32" type="uint" endian="big" value="0x51EDA7A0"/>
   <read var="payload" octets="{$size - 20}" type="raw"/>
   <read var="checksum" bits="32" type="uint" endian="big"/>
   <read var="size2" bits="32" type="uint" endian="big" value="{$size}"/>
  </loop>
 </decoder>
 <tag id="sie:xml_metadata" group="0" format="text/xml"/>

 <!-- SIE index block decoder:  v0=offset, v1=group -->
 <decoder id="1">
  <loop>
   <read var="v0" bits="64" type="uint" endian="big"/>
   <read var="v1" bits="32" type="uint" endian="big"/>
   <sample/>
  </loop>
 </decoder>
 <tag id="sie:block_index" group="1" decoder="1"/>

<!-- Stream-specific definitions begin here: -->
\end{lstlisting}

Notice this contains a description of the block structure described in
section~\ref{block-structure-sec}, and a tag \xml{sie:xml\_metadata},
the value of which is the contents of group 0.  It also contains a tag
and description of the format of the block indexes, described below.

The purpose of including this preamble is not to expect that a reader
program will extract how to read the SIE block structure from this
description; a sufficiently clever implementation could, but that
value of this is questionable.  Rather, it is to realize the design
goal that the format be completely self-describing.  Even without this
document, looking at an SIE file in a text editor will yield a good
chance of understanding the basic format.

To make possible some of the properties of the SIE format, the XML
representation used has some unusual features.  To support streaming
data, the XML format supports effectively overriding and modifying
metadata from earlier in the stream.  For example, the following two
samples are equivalent:

\begin{multicols}{2}
\begin{lstlisting}[language=siexml]
<ch id="42" name="test">
  <tag id="core:description">testing</tag>
</ch>
<ch id="42">
  <tag id="core:output_samples">74088</tag>
</ch>
\end{lstlisting}

\begin{lstlisting}[language=siexml]
<ch id="42" name="test">
  <tag id="core:description">testing</tag>
  <tag id="core:output_samples">74088</tag>
</ch>
\end{lstlisting}
\end{multicols}

This is an example of \emph{merging} behavior.  This can be used to
add information once it is known --- in the example above, the number
of output samples is clearly not known until all samples are written
to the file.

Some other metadata are atomic and \emph{replace} instead, such as
\xml{tag}s:

\begin{multicols}{2}
\vbox{
\begin{lstlisting}[language=siexml]
<ch id="42" name="test">
  <tag id="core:description">testing</tag>
</ch>
<ch id="42">
  <tag id="core:description">overridden</tag>
</ch>
\end{lstlisting}
}\vbox{
\begin{lstlisting}[language=siexml]
<ch id="42" name="test">
  <tag id="core:description">overridden</tag>
</ch>
\end{lstlisting}
}
\end{multicols}

Table~\ref{xml-merge} details the merging and replace behavior for
interpreting the SIE XML metadata.  \emph{Behavior} is either
\emph{merge} or \emph{replace}, the behaviors of which are described
above.  \emph{Key} is the attribute which is used to select which
existing element to merge with or replace.  If the key is
\emph{unique}, only one element of that type is allowed within its
parent.

\begin{table}
  \begin{center}
    \begin{tabular}{lll}
      \textbf{Element} & \textbf{Behavior} & \textbf{Key}  \\ \hline
      \xml{ch}         & merge             & \xml{id}      \\
      \xml{dim}        & merge             & \xml{index}   \\
      \xml{test}       & merge             & \xml{id}      \\
      \xml{data}       & replace           & \emph{unique} \\
      \xml{tag}        & replace           & \xml{id}      \\
      \xml{xform}      & replace           & \emph{unique}
    \end{tabular}
  \end{center}
  \caption{SIE XML merge and replace behavior}  \label{xml-merge}
\end{table}

\label{path-spec}
Additionally, to reduce metadata size and require less state from the
writer, there is a shortcut method to describe elements nested within
others.  The attribute names \xml{test}, \xml{ch}, and \xml{dim} are
reserved for all elements, and produce hierarchy as follows:

\begin{multicols}{2}
\vbox{
\begin{lstlisting}[language=siexml]
<tag 
   ch="2" dim="3" test="1"
   id="core:description">test</tag>
\end{lstlisting}
}\vbox{
\begin{lstlisting}[language=siexml]
<test id="1">
  <ch id="2">
    <dim index="3">
      <tag id="core:description">test</tag>
    </dim>
  </ch>
</test>
\end{lstlisting}
}
\end{multicols}

Note that to allow this to be implemented by a simple expansion, there
cannot be redundant element specifiers.  For example, the following
would be in error, as it could create two nested \xml{test} tags:

\begin{lstlisting}[language=siexml]
<test id="1">
  <tag 
     ch="2" dim="3" test="1"
     id="core:description">test</tag>
</test>
\end{lstlisting}

\label{inherit}
Finally, the \xml{ch} and \xml{test} elements support
\emph{inheritance}.  To inherit, set the \xml{base} attribute to the
\xml{id} of the element you wish to inherit from.  For instance, the
following are equivalent:

\begin{multicols}{2}
\vbox{
\begin{lstlisting}[language=siexml]
<ch id="2" name="old">
  <tag id="core:description">test</tag>
</ch>
<ch id="42" base="2" name="new"/>
\end{lstlisting}
}\vbox{
\begin{lstlisting}[language=siexml]
<ch id="2" name="old">
  <tag id="core:description">test</tag>
</ch>
<ch id="42" name="new">
  <tag id="core:description">test</tag>
</ch>
\end{lstlisting}
}
\end{multicols}

Inheritance is useful for creating many similar channels without
excessive duplication.  
%{\small FIXME The complete definition of a
%  \xml{ch} is considered to be a clone of the \xml{ch} identified by
%  \xml{base} (recursively interpreted through its \xml{base} if any),
%  with conflicting fields in the clone replaced by those specified in
%  the top level definition.}

\subsubsection{Index blocks (group 1)}

Index blocks, if present, allow an SIE file to be opened more quickly.
The payload of an index block consists of a sequence of 64-bit offset
and 32-bit groups; each offset and group correspond to one of the
blocks that are being indexed.  Index blocks appear \emph{after} the
blocks which they are indexing, allowing the SIE file to be streamed.

Each index block must cover a consecutive chunk of blocks ending with
the last block before the index block itself.  Index blocks that are
present in another later index block will be passed over to allow
reindexing of a file.

\subsection{XML metadata grammar}

\subsubsection{\xml{sie} element}

The \xml{sie} element is the top-level element of the SIE XML
metadata.  The \xml{version} attribute specifies the version of
SIE --- version \siever{} is described in this document.  The \xml{sie}
element can contain \xml{tag}s, \xml{decoder}s, channels (\xml{ch}),
and \xml{test}s.

The \xml{sie} element is described below in the \relaxng{} schema
description language.  Similar \relaxng{} syntax will follow all
other element descriptions below.

\begin{grammar}
SIE = element sie {
  attribute version { text },
  ( TopTag | Decoder | Channel | Test )*
}
\end{grammar}

Note that as an SIE stream is always capable of being appended to, the
\xml{sie} element should never be explicitly closed.  For parsing
purposes, consider ``\lit{</sie>}'' added to the end of the XML
metadata.

\subsubsection{\xml{tag} element}

The \xml{tag} element is the generic element for application-level
annotation, which allows arbitrary text as the tag name, and
completely arbitrary values.

\xml{tag} values are stored either directly in the XML, properly escaped:
\begin{lstlisting}[language=siexml]
<tag id="core:sample_rate">500</tag>
\end{lstlisting}
or as the contents of a group:
\begin{lstlisting}[language=siexml]
<tag id="setup_photo.1" group="27" format="image/jpeg"/>
\end{lstlisting}

For the \emph{group} form, the value of the tag is considered to be
the sequential concatenation of the payload of every block with the
specified group ID.

The \lstinline[language=siexml]!<tag group="42" decoder="17"/>! form
allows binary structures specifiable through the \xml{decoder}
mechanism. Extant SIE reader implementations are not capable of
applying \xml{decoder}s to \xml{tag} values yet, however.

\begin{grammar}
TestSpecs = ( ChannelSpec | ( ChannelSpec & DimSpec ) )
\end{grammar}
\begin{grammar}
TopTag = element tag {
  ( attribute test { UINT } | TestSpecs )?,
  TagBase
}
\end{grammar}
\begin{grammar}
TestTag = element tag {
  TestSpecs?,
  TagBase
}
\end{grammar}
\begin{grammar}
ChannelTag = element tag {
  DimSpec?,
  TagBase
}
\end{grammar}
\begin{grammar}
Tag = element tag {
  TagBase
}
\end{grammar}
\begin{grammar}
TagBase = (
  attribute id { text },
  ((attribute group { UINT }, DecoderOrFormat?) | text)
)
\end{grammar}
\begin{grammar}
DecoderOrFormat = ( attribute decoder { UINT } | Format )
\end{grammar}
\begin{grammar}
UINT = xsd:nonNegativeInteger
\end{grammar}

\subsubsection{\xml{decoder} element}

Each \xml{decoder} contains the machinery to decode the binary payload
from a particular data block. The same decoder may be used by any
number of different channels; the assignment of a particular decoder
ID to a group is mediated through the \xml{group} and \xml{ch}
directives. 

\begin{grammar}
Decoder = element decoder {
  Id,
  DecodeOps*
}
\end{grammar}
\begin{grammar}
DecodeOps = ( If | Loop | Read | Sample | Seek | Set )
\end{grammar}

\paragraph{Variables}

Decoder variables can contain either numbers or arbitrary binary data
of unbounded size.  Variables are always initialized to zero upon
entering a decoder.  The special variables \xml{v0}, \xml{v1}, \ldots,
\xml{v}$n$ are used by the \xml{sample} operator to compose output
result vectors.

\paragraph{Expressions}

Expressions start with a \xml{\{} and end with \xml{\}}.  Simple
arithmetic expressions are supported, and variable dereferencing via
the \xml{\$} prefix, i.e. the value of \xml{\{\$foo + 24\}} is 24 plus
the value of the variable \xml{foo}.

\begin{grammar}
EXPR = text  # either a literal number or an {...} expression
\end{grammar}

\paragraph{\xml{if} operator}

\xml{if} is used for conditional execution of code.  If the expression
in the \xml{condition} attribute evaluates to non-zero, the contents
of the \xml{if} are executed.

\begin{grammar}
If = element if {
  attribute condition { EXPR },
  DecodeOps*
}
\end{grammar}

\paragraph{\xml{loop} operator}

\xml{loop} is a general iteration mechanism.  By default, it repeats
its contents forever or until the decoder aborts upon a failed
\xml{read} operation.  If the optional \xml{var} attribute is present,
the numeric variable defined within (the \emph{loop variable}) will
incremented by one at the end of each iteration.

If the \xml{var} attribute is present, the optional \xml{start},
\xml{end}, and \xml{increment} attributes can be specified, each
containing an expression.

The \xml{increment} defaults to one; the loop variable will be
incremented by this value at the end of each loop iteration.  If this
is defined as an expression, the increment will be computed each
iteration.

The \xml{start} attribute is the value that the variable in the
\xml{var} attribute will begin the loop with.
\lstinline[language=siexml]!<loop var="foo" start="42">!
is exactly equivalent to
\lstinline[language=siexml]!<set var="foo" value="42"/><loop var="foo">!.

The \xml{end} attribute sets the termination condition; the loop runs
while $\xml{var} < \xml{end}$ when the increment is zero or positive,
and $\xml{var} > \xml{end}$ when the increment is negative.  If the
increment is variable, the termination condition style is computed
each time through the loop appropriately.

\begin{grammar}
Loop = element loop {
  (attribute var { text }, LoopOpts)?,
  DecodeOps*
}
\end{grammar}
\begin{grammar}
LoopOpts = (
  attribute increment { EXPR }?,
  attribute start { EXPR }?,
  attribute end { EXPR }?
)
\end{grammar}

\paragraph{\xml{read} operator}

The \xml{read} operator reads $n$ bits of payload from the current
position, where $n$ is either the value of the \xml{bits} attribute,
or eight times the value of the \xml{octets} attribute.  \xml{bits},
if used, must be a multiple of eight.  If neither size attribute is
present, all remaining data in the payload is read.  If there is less
data remaining than requested, the decoder terminates; this is the
primary means of termination for most decoders.

If the \xml{var} attribute is present, the named variable is set to
the value read, interpreted according to any \xml{type} and
\xml{endian} attributes that may be present.  If no \xml{type} or
\xml{endian} attributes are set, or if the \xml{type} attribute is set
to ``raw'', the data are assigned as a raw binary string.

The optional \xml{value} attribute asserts that the value read is
equal to the value of the attribute.  An error will result if this is
not the case.

The \xml{type} attribute has several possible values:  \\
\begin{tabular}{r@{: }l}
  ``int''   & Two's complement signed integer of 8, 16, 
              32, or 64 bits. \\
  ``uint''  & Unsigned integer of 8, 16, 32, or 64 bits. \\
  ``float'' & IEEE-754 floating point value of either 
              32 or 64 bits. \\
  ``raw''   & No interpretation; store unmodified binary buffer.
              The \xml{endian} attribute has no effect here.
\end{tabular}

The \xml{endian} attribute specifies the byte order to be used for
interpreting numerical data. The number \lit{0x00112233} would be
stored in the order \lit{0x00}, \lit{0x11}, \lit{0x22}, \lit{0x33} in
\xml{endian="big"} (also known as network or Motorola order), and as
\lit{0x33}, \lit{0x22}, \lit{0x11}, \lit{0x00} in
\xml{endian="little"} (Intel, Vax, x86 order).

\begin{grammar}
Read = element read {
  attribute var { text }?,
  ( attribute bits { EXPR } | attribute octets { EXPR }),
  ReadType,
  attribute value { EXPR }?
}
\end{grammar}
\begin{grammar}
ReadType = (
  Format
  | attribute type { "raw" }
  | ( attribute type { "int" | "uint" | "float" }
      & attribute endian { "big" | "little" } )
)
\end{grammar}

\paragraph{\xml{sample} operator}

\xml{sample} causes a data sample vector $[\xml{v0}, \xml{v1}, \ldots,
  \xml{v}n]$ to be added to the decoder's output queue.  All variables
of the form v0, v1, v2, etc.\ that are used in the decoder (through
\xml{loop}, \xml{set}, or \xml{read} operators) are included in the
output vector.

\begin{grammar}
Sample = element sample { empty }
\end{grammar}

\paragraph{\xml{seek} operator}

The \xml{seek} operator moves the location from which the next
\xml{read} operator will read data.  The read location is moved to the
position described in the attributes \xml{from} and \xml{offset}.
\xml{from} can either be ``start'', the start of the data;
``current'', the current position; or ``end'', the end of the data.
\xml{offset} is an expression evaluating to the number of octets
offset from the \xml{from} location.

\begin{grammar}
Seek = element seek {
  attribute from { "start" | "current" | "end" },
  attribute offset { EXPR }
}
\end{grammar}

\paragraph{\xml{set} operator}

\xml{set} is used to set the variable specified by the attribute
\xml{var} to the value of the attribute \xml{value}.  \xml{value} can
be either a literal number or an expression.

\begin{grammar}
Set = element set {
  attribute var { text },
  attribute value { EXPR }
}
\end{grammar}

\subsubsection{\xml{ch} element}

\xml{ch} is the container used to define engineering level data
channels.  It contains \xml{dim} elements whose
\xml{xform}, \xml{data}, and \xml{tag} elements have all the
information necessary to convert the raw decoder data vectors into
engineering unit data.  The optional \xml{name} attribute defines the
channel name.  All other information is contained in \xml{tag}
elements either directly under \xml{ch} or under a \xml{dim}
child.

To minimize redundancy, \xml{ch} supports an inheritance mechanism
through its \xml{base} attribute.  See page~\pageref{inherit} for more
details.

The optional \xml{group} attribute specifies a single group to be
associated with the channel.  This will be most often seen using
inheritance inside a test definition:
\begin{lstlisting}[language=siexml]
<test id="1">
  <ch id="9" base="0" group="5" name="th1@rpm.RN_2"/>
  <ch id="10" base="1" group="5" name="th1@coolant.RN_2"/>
</test>
\end{lstlisting}

%% More than one group can contribute to a channel:  the \xml{dim}
%% \xml{group} attribute overrides any \xml{group} attribute of the
%% parent channel.

\begin{grammar}
Channel = element ch {
  Id,
  attribute base { UINT }?,
  attribute name { text }?,
  attribute group { UINT }?,
  Private?,
  (Dimension* & ChannelTransform* & ChannelTag*)
}
\end{grammar}

\subsubsection{\xml{dim} element}

The \xml{dim} element contains the definition for axis \xml{index}.
The first axis is index zero.  Decoder data for the dimension comes
from the \xml{dim}'s \xml{group} attribute, or if that does not exist,
the enclosing \xml{ch}'s \xml{group} attribute.

\begin{grammar}
Dimension = element dim {
  attribute index { UINT },
  attribute group { UINT }?,
  ( Transform? & Data? & Tag* )
}
\end{grammar}

\subsubsection{\xml{xform} element}

\xml{xform} defines a linear tranform using the \xml{scale} (i.e. slope)
and \xml{offset}" (i.e. intercept) attributes.  An indexed mapping
transform (essentially an array lookup) uses \xml{index\_ch} and
\xml{index\_dim} to define the output.
%%  Currently dimension 0 of the
%% specified channel is used implicitly as the "from" part of the
%% map. <= KLUDGE check this statement

\begin{grammar}
Transform = element xform { Transforms }
\end{grammar}
\begin{grammar}
TestTransform = element xform { ChannelDimSpec, Transforms }
\end{grammar}
\begin{grammar}
ChannelTransform = element xform { DimSpec, Transforms }
\end{grammar}
\begin{grammar}
Transforms = (LinearTransform | IndexTransform)
\end{grammar}
\begin{grammar}
LinearTransform = (
  attribute scale { REAL },
  attribute offset { REAL }
)
\end{grammar}
\begin{grammar}
IndexTransform = (
  attribute index_ch { UINT },
  attribute index_dim { UINT }
)
\end{grammar}

\subsubsection{\xml{data} element}

\lstinline[language=siexml]!<data decoder="d" v="v">! is used to
define the data for a particular \xml{dimension} as the vector element
$v$ from decoder $d$.  Thus
\begin{lstlisting}[language=siexml]
<dim index="2">
  <data decoder="7" v="3">
</dim>
\end{lstlisting}
assigns element v3 of decoder 7's output to dimension 2 of the
channel.

\begin{grammar}
Data = element data {
  DecoderOrFormat,
  attribute v { UINT }
}
\end{grammar}

\subsection{Data rendering algorithm}

The first step in getting the data model representation of an SIE
channel's data is to look at the \xml{ch} element defining the
channel's metadata.

The following information must be obtained for each dimension $d$:
which group the data are going to be read from ($g$), which decoder
will be used to decode the data ($D$), which \xml{v} of the decoder
will be assigned to the dimension ($v_d$), and optionally, what
transform will be applied to the output of the decoder ($x_d$).  Note
that for this SIE version, the group and decoder for all dimensions
must be the same, as there is no way to specify how multiple disparate
decoder outputs be joined.  If any of the required information above
is missing, the channel is \emph{abstract} and does not have data
associated with it.  This is common in the case of base channels which
are only used to hold common metadata for other channels.  Usually,
these channels will have the \xml{private} attribute set so the reader
knows not to try to access them.

To determine the group, look at the \xml{group} attribute for each of
the \xml{dim} elements, \emph{or} the \xml{group} attribute of the
\xml{ch} element.  The decoder ID is defined by the \xml{decoder}
attribute in each dimension's \xml{data} element, and the decoder
\xml{v} by the \xml{v} attribute.  Finally, the transform is defined
by the dimension's \xml{xform} element, if any.

Knowing $g$, $D$, all $v_d$, and all $x_d$, the algorithm to render
the channel data is relatively straightforward:

\begin{tabbing}
For \= each block $b$ in the SIE stream with group ID $g$: \\
    \> Run decoder $D$ on the payload of $b$, producing multiple
       output vectors. \\
    \> For \= each decoder output vector $V$: \\
    \>     \> Start building up an output row $r$. \\
    \>     \> For \= each dimension $d$ in the channel: \\
    \>     \>     \> Let $o$ be the decoder output $V[v_d]$. \\
    \>     \>     \> If the transform $x_d$ exists, apply it to $o$. \\
    \>     \>     \> Assign $r[d]$ to be the value of $o$. \\
    \>     \> Append the row $r$ to the channel data.
\end{tabbing}

The result of this algorithm will be the channel data in the form of
the universal SIE data model as described in
section~\ref{data-model-sec}.

\section{Implementation}

\subsection{Writing SIE}

The design of SIE is biased to be simple and efficient to write.  An
SIE writing implementation can output almost any binary format
desired; the binary data need only be packaged in self-contained
format in SIE blocks and have the appropriate SIE XML metadata
included defining how to interpret it.

As an example, consider an extremely simple data logger.  This
theoretical data logger collects a single channel of time series data,
and stores the data in 16-bit signed integers in big-endian byte
ordering to on-board flash.  Furthermore, the logger only has 128
bytes of usable RAM.

Outputting SIE from this device is quite feasible.  Pre-computed XML
metadata segments for this device's single output format are stored in
ROM.  In between these segments are ``holes'' where computed metadata
values needs to go.

To output the data in SIE format, the first XML section (containing
from the beginning of the SIE XML preamble up to the first variable
metadata) is output in a group 0 block.  While this is larger than the
RAM available, it can be done incrementally since the size is known
ahead of time.  Then, the first piece of variable metadata is computed
and a group 0 block containing just that metadata is output.  This
process of outputting pre-computed XML and outputting computed
metadata values continues until the end of the XML is reached.  Then,
the size of the data is computed, a block header for the data is
output, the actual data is output, and a block footer for the data is
output.

All CRCs can be written as zero, though since the CRCs for the
constant XML sections could be precomputed they might as well be
written out.  Index blocks are also optional and would not be written
in this scenario.

While this is an extreme example, it demonstrates that writing a
fully-compliant SIE file needs little more than ROM space to store the
XML.

\subsubsection{Writing guidelines}

%% \incomplete{}

\begin{itemize}
\item Choose data representations that are self-contained (i.e., in
  \xml{somat:histogram}, each row is a complete specification of a
  bin, containing both its count and its position in engineering
  units).
\item Choose block sizes that strike a balance between overhead (as
  there are 20 bytes of overhead for each block written; 32 and some
  small fraction with indexes) and seekability (the block is the
  natural unit of reading granularity in an SIE file; don't write out
  blocks larger than the RAM of the reading computer).
\item If at all possible, use CRCs on XML and index blocks even if in
  general CRCs cannot be used on all data due to performance
  limitations.
\end{itemize}

%% \subsubsection{The eDAQ SIE writing implementation}

%% \incomplete{}

\subsection{Reading SIE}

On an abstract level, reading an SIE file is relatively easy --- as
described earlier, the format was designed to be self-documenting
enough that looking at an SIE file in a text editor should reveal
enough information for a human to be able to manually extract
engineering data from the file.  Writing a computer program to read a
\emph{specific} SIE file is also relatively straightforward.  However,
writing a computer program to read \emph{all} SIE files is rather
complicated.

Since a design goal of SIE was to allow simple and flexible writing
implementations, a lot of the complexity is shifted onto the reading
side.  The nature of the SIE format basically mandates than a general
reader implementation must contain at least an XML parser, a parser
for the decoder expression language, a decoder language interpreter or
compiler, a robust implementation of the SIE data rendering algorithm,
and an API general enough to deal with the wide variety of data
representable in an SIE file.  This is a major software undertaking.
Additionally, the presence of limited or faulty SIE reading
implementations weakens the advantages of the format --- SIE writers
will not be free to use the inherent flexibility of the format to
their advantage if the reading implementations used cannot handle the
full specification.

The authors, aware of this problem, sought to write a single general
SIE reader implementation, \libsie{},  and release it under a liberal
freely-available software license, so as to be available to all.
%%   This
%% software, \libsie{}, is described below.

%% \subsubsection{The \libsie{} reading implementation}

%% \section{Future work}

%% \subsection{Support for streaming data}

%% \subsection{Corrupt data recovery}

%% \section{Conclusions}

%% \appendix

%% \section{SIE RelaxNG XML specification}

%% \verbatiminput{../../../sie/SIE.rnc}

\appendix

\section{License}

\begin{verbatim}
Copyright (C) 2005-2010  HBM Inc., SoMat Products

This document is free software; you can redistribute it and/or
modify it under the terms of version 2.1 of the GNU Lesser General
Public License as published by the Free Software Foundation.

This document is distributed in the hope that it will be useful, but
WITHOUT ANY WARRANTY; without even the implied warranty of
MERCHANTABILITY or FITNESS FOR A PARTICULAR PURPOSE.  See the GNU
Lesser General Public License for more details.

You should have received a copy of the GNU Lesser General Public
License along with this document; if not, write to the Free Software
Foundation, Inc., 51 Franklin Street, Fifth Floor, Boston, MA
02110-1301 USA
\end{verbatim}

\bibliography{bibliography}{}
\bibliographystyle{plainnat}

\end{document}
